%% APPENDIX D -- the false-alarm expectation, the control ensemble, and the
%% execution-block quality criterion.  Owner: r14-chance.  Carries
%% \label{app:falsealarm} (04_method.tex, 08_backmatter.tex) and
%% \label{sec:blockquality} (04_method.tex).
\section{The false-alarm expectation and the control ensemble}
\label{app:falsealarm}

This appendix gives the chance expectation for the unattributed crossings,
what it is conditioned on, and what the control ensemble calibrating it can
support.

\subsection{What the expectation is conditioned on}

\emph{One attribution set serves both sides of the comparison}, the set of
\S\ref{sec:technosearch}: the mask of Appendix~\ref{app:linemask} within
$\pm\MaskHalfKms$\,km\,s$^{-1}$ in the star's own rest frame, sulphur monoxide
excluded. So the \MkNSODecide{} sulphur monoxide coincidences are unattributed
crossings, and a control maximum on one of those transitions is an
unattributed exceedance. A systemic velocity is measured for the stars of
\CeNVsys{} of the census windows and the others are taken at rest, which moves
the masked share by \CeVsysSensPct{} points.

The expectation is built from the control positions, the only empirical null
the data provide, and each Class~A window contributes two quantities
measured on itself. The fraction of its own \NCtrl{} control positions that
reaches the trigger is the probability that a position in that field does so
--- no cell count enters it and no form is assumed for the tail --- and
summed over the windows \CeExpTrig{} are expected to trigger by chance. The
share of its own channel grid lying inside that same set is the probability
that such a position would be attributed, a control maximum being free to
fall on any channel. Weighted by
the first that share is \CeMaskPct{} per cent, against a median of
\CeMaskPctMed{} per cent per window, because \CeNFullMask{} narrow
line-tuned windows lie wholly inside the mask and their fields, which hold
real emission, are where a control position most often triggers. What is
left is \CeExp{} expected unattributed Class~A crossings against \CeObs{}
observed.

Both sides run over the \CeNWin{} Class~A windows of the primary census, so
the expectation is a sum over rows and nothing is carried at an assumed rate.
A further \CeNWinTail{} Class~A windows arrived with the archival tail and lie
outside it, together with the \CeNTailCrossWord{} unattributed crossings
toward \CeTailStar{} they hold, which are reported beside the comparison and
enter neither side of it. No window holds more than one primary crossing, so
both sides count windows; the weaker cells behind those maxima are counted at
their own grain in \S\ref{sec:technosearch}.

The one assumption the searched data cannot test is that a star behaves like
a control position in its own field, and the reserved blocks measure it in
the quantity used here. Their stars reach the trigger in \StHoObs{} of
\StHoNFine{} fine-channel windows where the control positions of those same
windows predict \StHoExpFine: a ratio of \StHoRatio, with \StHoLo--\StHoHi{}
allowed by resampling the reserved blocks, and no two of the events in one
block. In the \StHoNCoarse{} coarse-channel reserved windows no star reaches
the trigger, against \StHoExpCoarse{} predicted. The same blocks place the
stellar rank below one half (Appendix~\ref{app:radius}), which is the same
departure read in the body of the distribution rather than in its tail and
runs the same way. Carrying the measured ratio through gives \CeExpX{}
against the \CeExp{} above, and the observed \CeObs{} sits inside the range
either way; the ratio is itself consistent with unity, so the uncorrected sum
is the one quoted.

The \CeLo--\CeHi{} quoted beside the expectation is what \CeNBoot{}
resamplings of the execution blocks allow, each window's rate redrawn from
its own \NCtrl{} probes: windows share blocks and are not independent
trials, and the ensemble is finite. Those two terms alone give
\CeLoNull--\CeHiNull, so what widens the range is the calibration above. An
added population of as many as \CeAddBlind{} events would be absorbed by that
width, and that is the limit on what the comparison can exclude.

Every Class~B crossing in the survey comes from the one execution block of
\S\ref{sec:blockquality}, and over both classes together \CeExpSurv{} are
expected against the same \CeObs{} observed: a coarse single pointing holds
too few cells for its own field to reach the trigger by chance, and only
\NCtrlWinObs{} of the \NWindows{} windows hold a control position above
$T_{\rm trig}=\StageNullTrig$ at all. The Class~A comparison is the survey's.

\emph{The statistic is always the one measured over the window's whole field.}
\StNRep{} windows admit integrations inside the primary beam that a narrower
field would discard, recovering a median \StRepOnsrc{} times the on-source
time and moving the statistic by a median \StRepDT; the whole-field value is
the one adopted, on \StRepAdopted{} rows of Table~\ref{tab:ledger}, and each
window's control ensemble is measured with it, so neither side of the
comparison mixes two fields. Where the extraction catalogue's disposition
column disagrees --- on \StRepFell{} crossings, \StRepAttrClause{} --- the
ledger is the authority.

\subsection{Execution-block quality}
\label{sec:blockquality}

The comparison with the control positions assumes each block's ensemble
behaves like noise, and one block makes that false: every one of its \NCtrl{}
probes exceeds the trigger, the smallest at \BqFailCtrlMin. A criterion for
the condition is therefore applied to every block rather than to the one that
prompted it.

\emph{Two criteria are built from that statistic and they are not the same
test.} The statistic is the median of a window's own \NCtrl-position control
ensemble, which separates a field with a source in it from a field that is no
null at all. A \emph{window} is flagged when that median exceeds \BqFactor{}
the median of the same quantity over the survey, \DqThresh{} against
\DqSurvMed; the flag is a property of the window and cuts across attribution
by design, firing on \DqNWin{} of the \NWindows{} windows in
\BkFlagNBlock{} blocks toward \DqNStar{} stars (\DqStars), and every one of
those windows carries a crossing, attributed and unattributed alike. A
\emph{block} is judged on the same median divided by the median over the
windows of its own class, since a maximum over a larger
channel$\times$drift grid sits higher and one threshold should mean the same
thing in both; a block fails when the median of that ratio over its windows
exceeds \BqFactor{} the scale of its class. The factor is the same in
both. Neither was fixed in advance: the condition
was noticed in one block and the criterion written afterwards, so it is post
hoc, and what makes it trustworthy is that it does not depend on where the
line was drawn --- the failing block sits at \BqFailQ{} and the next worst at
\BqNextQ, so any threshold between them selects the same single block.

Applied to all \BqNBlocks{} blocks of the primary census, where the median
block sits at \BqQMed, the criterion fails \BqNFail: \BqFailBlock{} toward
\BqFailStar. Nothing else is close, and the unscaled form of the rule fails
the same \BqRawNFail{} block. Its \BqFailNWin{} windows are all coarse-channel and
hold the survey's entire Class~B crossing population, so it carries
\EvNCrossFlag{} of the \EvNCrossRaw{} threshold crossings --- \EvNRankFlag{}
of them outranking all \NCtrl{} control positions of their own window --- and
no Class~A statistic here has ever contained it. Those \EvNCrossFlag{} appear
in Table~\ref{tab:ledger} with their block named and enter no count
elsewhere.

\subsection{The control ensemble: its calibration, and its resolution}

The ring geometry is set in \S\ref{sec:statistic}; its calibration and its
resolution are the subject here. Over the \VerNTested{} census windows that
retain a control vector the stellar ranks are consistent with the one half
exchangeability implies --- median \BkRankBalSurv{} with each star carrying
equal weight, \BkRankBalSurvLo{} to \BkRankBalSurvHi{} --- in both array
strata, the compact-array windows measured on each observation's own antenna
table rather than on a 12\,m beam. That is a weak statement and not a strong
one, since a shift smaller than about \VerKsPowerAca{} could not have been
seen in a sample of this size. The reserved blocks lean further from one
half, and Appendix~\ref{sec:heldout} compares the two samples. The
expectations above are summed from the exceedance rate measured at each
window's control positions, so they do rest on that transfer and carry its
measured accuracy.

\emph{The rank is not a significance test, and the reason is structural.}
Ranking the star against \NCtrl{} control positions of the same field
decides which windows deserve examination, but $1/(\NCtrl+1)=\RankFloorVal$
is the \emph{resolution} of the ensemble --- the smallest rank it can
express --- and not the probability that noise produced the window. A survey
of \NWindows{} windows would need a rank \BonfRankRatio{} times finer, and
correlation makes it coarser still: controls inside the synthesised beam are
not independent, so an annulus resolves no better than
\RankEffLo--\RankEffHi{} however many probes are drawn into it, with the
primary beam bounding how many a single pointing contains. That floor is the
design's central limitation: no single window here can be significant on the
rank alone, which is why confirmation requires independent recurrence.

\subsection{The trials budget}

The searched cells counted in \S\ref{sec:summary} are not independent
trials: channel smoothing and a drift grid stepped twice per channel
correlate neighbouring ones, which alone would leave \StGridFrac{} of them
independent. The control ensembles say otherwise, their maxima sitting at
\CtrlMaxRatioMed{} of the prediction for the maximum of the full searched
count treated as independent and Gaussian --- a budget predicting
\NCtrlWinExp{} windows with a control above $T_{\rm trig}$ against the
\NCtrlWinObs{} observed --- while fitting each window's effective count
against its own ensemble returns \StNindFrac{} of its searched cells. Nothing
turns on which is right, because the expectation above is summed from measured
rates and not from a cell count: by the cell-count route the number reaching
the trigger is \StExpModel{} rather than \CeExpTrig, \StModelAgreePct{} per
cent away.

\emph{The on-star crossings do exceed that budget, and the attribution
accounts for it.} A Gaussian tail over every searched cell predicts
\CellExpect{} stellar exceedances against the \NCross{} found, but \MkAttr{}
are attributed by the set above, \EvNCrossFlag{} come from the block of
\S\ref{sec:blockquality} and \CeNTailCross{} fall outside the census windows,
leaving \CeObs{} against the \CeExpSurv{} predicted. No excess survives the
attribution.

\emph{One stratum departs from uniformity, and the departure is where the
line emission is.} Counting the windows whose stellar peak both reaches the
trigger and exceeds all \NCtrl{} control positions of its own window, the
$\NFine$ fine-channel windows hold \NStageOneFine{} against
\ExpStageOneFine{} expected, while the coarse stratum, three times larger
and predicting \ExpStageOneCoarse, contains \NStageOneCoarse. A coarse
channel dilutes a narrow circumstellar line, so the class that can resolve
one is the class that finds it --- and an unresolved carrier would appear in
Class~A first for the same reason. Reaching the trigger is holding a
crossing, so each of those windows holds one and \CeFirstAttr{} hold a
crossing the mask attributes: the departure is a statement about the windows
carrying line emission and not about the fine stratum at large. Drop the
trigger condition and the star leads its own ring in \BkFirstBelow{} further
windows, at \BkFirstBelowTLo{} and \BkFirstBelowTHi, giving \NsrFirstObs{}
over all \NWindows{} windows against \NsrFirstExp{} expected. One quantity,
two conditions; the chance comparison uses the conditioned one.

\emph{What is left is bounded, and the bounds are null.} The real stellar
ranks agree at $\pv{\StarPseudoP}$ with the \NPseudo{} ranks each control
takes against the other probes of its own window; radial suppression is below
\RingRhoBoundPct{} per cent in the mean; and scrambling a window so nothing
coherent survives along a drift track carries \LNNWindows{} past the rank
floor and changes no disposition.

%% The chance expectation here is summed from the measured control exceedance
%% rates and each window's own masked grid share (chance_v414.py, round 310),
%% so no cell count and no Gaussian tail enters it and nothing in it is
%% conditioned on the crossings; the cell-count route is a cross-check only,
%% and the rank-conditioned expectation (\ClustStageMean against
%% \ClustStageNObs) is a pipeline diagnostic that must not be reachable as
%% "the" expectation.  tab:ctrlsep, tab:excess, tab:arraysplit and
%% fig:ctrldiag are generated and nothing cites them; r15-deposit, 2026-10-08:
%% they used to be described as "in the data release", and there is no data
%% release, so they are simply unpublished diagnostics.
%% tab:ledger lives here rather than in Sec. 5, beside the false-alarm
%% accounting that explains how many of its rows are expected.
\begin{table*}
\centering
\caption{\textbf{The crossing ledger, from which every disposition in this
paper can be reproduced.} All \EvNCrossRaw{} threshold crossings,
including the \EvNCrossFlag{} in the failing block of \S\ref{sec:bdblock}.
Every row is the primary crossing of its window; the \CkNSecondary{} weaker
cells behind those maxima are not listed individually, and the
\CkNSecUnattrWord{} of them the mask does not attribute are described in
\S\ref{sec:technosearch}. The columns are
the star, its execution block, the ALMA band, the crossing
frequency as the correlator delivered it, or, where no peak frequency was
measured, recovered from the line offset, and the same frequency in the star's
rest frame, the search statistic $T_\star$, the largest statistic
$T_{\rm ctrl}$ among the window's \NCtrl{} control positions, the nearest
masked transition and the stellar-frame offset $\Delta v_\star$ from it, then
the recurrence test: $n_{\rm rep}$ covering repeat blocks spanning $\Delta t$
days, the statistic $T_{\rm pers}$ a carrier persisting at the discovery flux
would have produced at the predicted stellar-frame cell, the statistic
$T_{\rm rep}$ measured there at the carried drift, and the exclusion of
Eq.~\ref{eq:excl}, which three of these columns reproduce, two of them also
giving the limit defined in Appendix~\ref{app:radius}. Where the emission is
present again in a covering block the test returns a recovery and not an
exclusion, so both $T_{\rm rep}$ and $\sigma_{\rm excl}$ are blank and the
disposition says which it was.
$T_\star$ and $T_{\rm ctrl}$ come from the same measurement of the
same window; $T_{\rm ctrl}$ gates nothing and says only which crossing was
examined first, and the drift-maximised reading of the cell is given in
Table~\ref{tab:cand} for the two leading events and is not tabulated for
the rest. Rows whose window is data-quality flagged
(\S\ref{sec:blockquality}) are marked, and on those $T_{\rm pers}$ and the
exclusion are withheld: both are built from the discovery amplitude, and in a
window whose control ring stands as high as its star that amplitude cannot be
assigned to a point source at the star. \BkFlagNCrossOther{} of those rows
lies outside the failing block --- \BkFlagOtherStar, foreground cloud
emission along its sightline --- and
is withheld for that reason and not for its block. A crossing is attributed
by the set defined above, and a dagger marks the \MkNSODecide{} rows whose
coinciding species that set excludes. The transition named is the nearest in
the star's frame, so where it is far from every transition it is the magnitude
and not the name that carries information. Block identifiers omit the constant
\texttt{A002\_} prefix.}
\label{tab:ledger}
\tiny
%% GENERATED by ledger_v410.py -- do not hand-edit.
%% One row per threshold crossing, one epoch convention, stellar-frame attribution.
%% part 1 of 2, 28 rows; the split is made in the generator so the float cannot outgrow its page.
\begin{tabular}{@{}l@{\hspace{0.9pt}}l@{\hspace{0.9pt}}c@{\hspace{0.9pt}}r@{\hspace{0.9pt}}r@{\hspace{0.9pt}}r@{\hspace{0.9pt}}r@{\hspace{0.9pt}}l@{\hspace{0.9pt}}r@{\hspace{0.9pt}}r@{\hspace{0.9pt}}r@{\hspace{0.9pt}}r@{\hspace{0.9pt}}r@{\hspace{0.9pt}}r@{\hspace{0.9pt}}l@{}}
\hline
Star & block & B & $\nu_{\rm topo}$ (GHz) & $\nu_{\star}$ (GHz) & $T_\star$ & $T_{\rm ctrl}$ & line & $\Delta v_{\star}$ & $n_{\rm rep}$ & $\Delta t$ (d) & $T_{\rm pers}$ & $T_{\rm rep}$ & $\sigma_{\rm excl}$ & disposition \\
\hline
BD+05 1668 & Xc7fa6f\_X28a8 & 6 & 240.1016 & 240.1054 & 73.398 & 74.69 & CS($5$--$4$) & -5911.9 & 29 & 0.9--113 & n/a & +0.15 & n/a & unattributed; flagged \\
BD+05 1668 & Xc7fa6f\_X28a8 & 6 & 224.7422 & 224.7458 & 48.303 & 83.74 & H$_2$CO($3_{1,2}$--$2_{1,1}$) & -1264.5 & 29 & 0.9--113 & n/a & +0.37 & n/a & unattributed; flagged \\
BD+05 1668 & Xc7fa6f\_X28a8 & 6 & 226.7422 & 226.7458 & 41.355 & 75.48 & CN($2_{3/2}$--$1_{1/2}$) & +88.0 & 29 & 0.9--113 & n/a & +0.66 & n/a & unattributed; flagged \\
BD+05 1668 & Xc7fa6f\_X28a8 & 6 & 242.1953 & 242.1992 & 39.586 & 83.29 & CS($5$--$4$) & -3349.2 & 29 & 0.9--113 & n/a & +0.77 & n/a & unattributed; flagged \\
$\beta$~Pic & Xf5d76d\_Xeb8 & 3 & 115.2603 & 115.2711 & 30.765 & 7.34 & CO($1$--$0$) & -0.2 & 4 & 0.0--53 & n/a & --- & n/a & attributed; recurs \\
$\beta$~Pic & Xf5d76d\_Xe19 & 3 & 115.2603 & 115.2711 & 27.684 & 6.48 & CO($1$--$0$) & -0.2 & 4 & 0.0--53 & n/a & --- & n/a & attributed; recurs \\
HD 48370 & Xc26103\_X155a & 6 & 230.5117 & 230.5243 & 26.834 & 18.36 & CO($2$--$1$) & -17.8 & --- & --- & n/a & --- & n/a & attributed; flagged \\
$\beta$~Pic & Xf4df6f\_Xc7 & 3 & 115.2620 & 115.2711 & 17.907 & 5.76 & CO($1$--$0$) & -0.2 & 4 & 53 & n/a & --- & n/a & attributed; recurs \\
$\beta$~Pic & Xf5d76d\_Xcc5 & 3 & 115.2604 & 115.2711 & 17.288 & 5.52 & CO($1$--$0$) & -0.2 & 4 & 0.0--53 & n/a & --- & n/a & attributed; recurs \\
$\beta$~Pic & Xf5d76d\_X32f1 & 3 & 115.2606 & 115.2713 & 14.654 & 7.27 & CO($1$--$0$) & +0.4 & 4 & 0.0--53 & n/a & --- & n/a & attributed; recurs \\
$\beta$~Pic & Xd9668b\_X3a90 & 6 & 230.5161 & 230.5378 & 11.681 & 9.12 & CO($2$--$1$) & -0.3 & 5 & 1.0--1983 & n/a & --- & n/a & attributed; recurs \\
$\beta$~Pic & X6f1341\_X1484 & 6 & 230.5284 & 230.5379 & 9.832 & 8.20 & CO($2$--$1$) & -0.1 & 5 & 25--1984 & n/a & --- & n/a & attributed; recurs \\
$\beta$~Pic & X7116f1\_X1785 & 6 & 230.5272 & 230.5384 & 9.677 & 14.13 & CO($2$--$1$) & +0.5 & 5 & 25--1959 & n/a & --- & n/a & attributed; recurs \\
$\beta$~Pic & Xd9668b\_Xa9df & 6 & 230.5160 & 230.5377 & 8.948 & 8.41 & CO($2$--$1$) & -0.4 & 5 & 1.0--1984 & n/a & --- & n/a & attributed; recurs \\
$\beta$~Pic & Xa7a216\_X23e4 & 6 & 230.5273 & 230.5377 & 7.988 & 6.03 & CO($2$--$1$) & -0.4 & 5 & 601--1313 & n/a & --- & n/a & attributed; recurs \\
61 Vir & Xc079b5\_X82f & 7 & 346.3138 & 346.3232 & 5.958 & 5.65 & SO($9_{8}$--$8_{7}$) & -177.6 & 3 & 0.1--3.0 & 9.7 & +0.39 & 4.9 & unattributed; no recurrence \\
HD 285968 & Xd98580\_X3bf6 & 6 & 230.5021 & 230.5447 & 5.956 & 8.12 & CO($2$--$1$) & +8.7 & 3 & 13--44 & n/a & --- & n/a & attributed; recurs \\
HD 285968 & Xdb7ab7\_X5b4e & 6 & 230.5116 & 230.5447 & 5.868 & 7.22 & CO($2$--$1$) & +8.7 & 3 & 5.0--44 & n/a & --- & n/a & attributed; recurs \\
HD 285968 & Xdb4b9a\_X108f & 6 & 230.5100 & 230.5447 & 5.835 & 7.94 & CO($2$--$1$) & +8.7 & 3 & 5.0--39 & n/a & --- & n/a & attributed; recurs \\
CP$-$72 2713 & Xff0235\_X4a6d & 7 & 344.2697 & 344.2965 & 5.809 & 5.68 & SO($8_{8}$--$7_{7}$) & -12.3$^{\dagger}$ & 5 & 0.1--2.1 & 6.3 & -0.82 & 4.8 & unattributed; no recurrence \\
HD~139084~B & Xcd8029\_Xb6b0 & 7 & 345.8241 & 345.8241 & 5.763 & 17.70 & CO($3$--$2$) & +24.4 & --- & --- & --- & --- & --- & attributed \\
HD 14055 & Xfff3d8\_Xd46 & 7 & 331.7699 & 331.7664 & 5.575 & 5.99 & $^{13}$CO($3$--$2$) & +1068.7 & 30 & 0.9--650 & 170.4 & -0.02 & 5.6 & unattributed; no recurrence \\
HD 23484 & X10b22d7\_X2445 & 6 & 230.7904 & 230.8006 & 5.520 & 5.55 & CO($2$--$1$) & +341.5 & 4 & 0.1--2.1 & 10.4 & +1.25 & 4.3 & unattributed; no recurrence \\
HD 161868 & Xff45a6\_X10ca8 & 7 & 345.6338 & 345.6491 & 5.481 & 6.16 & SO($2_{3}$--$2_{1}$) & -48.0$^{\dagger}$ & 13 & 0.0--6.1 & 92.2 & +0.24 & 5.5 & unattributed; no recurrence \\
HIP 10679 & Xa95c04\_X14cd & 6 & 220.4185 & 220.4035 & 5.456 & 6.76 & $^{13}$CO($2$--$1$) & +6.5 & --- & --- & --- & --- & --- & attributed \\
HD~14055 & X11adad7\_X25c5e & 7 & 345.4564 & 345.4303 & 5.448 & 6.82 & SO($2_{3}$--$2_{1}$) & -237.8 & 30 & 0.1--653 & 14.0 & +0.03 & 5.1 & unattributed; no recurrence \\
HD 207129 & Xc19d6f\_X5706 & 6 & 230.4226 & 230.4070 & 5.433 & 5.52 & CO($2$--$1$) & -170.3 & 4 & 0.8--425 & 69.7 & -0.93 & 5.5 & unattributed; no recurrence \\
$\beta$~Pic & Xd9668b\_X3a90 & 6 & 241.7512 & 241.7739 & 5.420 & 5.81 & CS($5$--$4$) & -3869.7 & 1 & 1.0 & 5.7 & +0.72 & 3.4 & unattributed; no recurrence \\
\hline
\end{tabular}

\end{table*}

%% A `table*` cannot break across pages, so `ledger_v410.py` writes the body
%% in parts of at most 28 rows and asserts it has enough floats for them.
%% This is the continuation; the first caption carries the definitions.
\begin{table*}
\centering
\caption{\emph{The crossing ledger, continued.} The remaining
\CxLedgerPartTwo{} of the \EvNCrossRaw{} threshold crossings, in the same
order and with the same columns as Table~\ref{tab:ledger}, whose caption
defines them. These too are primary crossings, one to a window.}
\label{tab:ledgercont}
\tiny
%% GENERATED by ledger_v410.py -- do not hand-edit.
%% One row per threshold crossing, one epoch convention, stellar-frame attribution.
%% part 2 of 2, 28 rows; the split is made in the generator so the float cannot outgrow its page.
\begin{tabular}{@{}l@{\hspace{0.9pt}}l@{\hspace{0.9pt}}c@{\hspace{0.9pt}}r@{\hspace{0.9pt}}r@{\hspace{0.9pt}}r@{\hspace{0.9pt}}r@{\hspace{0.9pt}}l@{\hspace{0.9pt}}r@{\hspace{0.9pt}}r@{\hspace{0.9pt}}r@{\hspace{0.9pt}}r@{\hspace{0.9pt}}r@{\hspace{0.9pt}}r@{\hspace{0.9pt}}l@{}}
\hline
Star & block & B & $\nu_{\rm topo}$ (GHz) & $\nu_{\star}$ (GHz) & $T_\star$ & $T_{\rm ctrl}$ & line & $\Delta v_{\star}$ & $n_{\rm rep}$ & $\Delta t$ (d) & $T_{\rm pers}$ & $T_{\rm rep}$ & $\sigma_{\rm excl}$ & disposition \\
\hline
HR~1010 & Xc6fa08\_X1053 & 6 & 230.5133 & 230.5287 & 5.410 & 5.65 & CO($2$--$1$) & -12.0 & 7 & 2.0--14 & 8.4 & $<$+4.60 & $>$2.1 & attributed; not recovered \\
ALMA J153702653-33192492 & Xff0235\_X8392 & 6 & 230.5222 & 230.5392 & 5.405 & 13.41 & CO($2$--$1$) & +1.6 & 6 & 18--141 & 6.2 & $<$+4.15 & $>$1.4 & attributed; not recovered \\
HD 110411 & Xe45e29\_Xb7d4 & 6 & 219.4538 & 219.4357 & 5.393 & 5.85 & C$^{18}$O($2$--$1$) & -170.2 & 1 & 10 & 6.7 & +0.62 & 3.8 & unattributed; no recurrence \\
HD 31392 & X129754b\_X6946 & 6 & 231.1998 & 231.2231 & 5.289 & 5.81 & H($30\alpha$) & -876.3 & 8 & 0.1--41 & 15.7 & -0.26 & 5.1 & unattributed; no recurrence \\
HD 14055 & Xff99e1\_X24be0 & 7 & 331.7078 & 331.7035 & 5.277 & 5.69 & $^{13}$CO($3$--$2$) & +1011.6 & 30 & 0.9--652 & 252.2 & +0.15 & 5.3 & unattributed; no recurrence \\
AU Mic & Xa42f75\_X319 & 6 & 230.6864 & 230.6705 & 5.274 & 6.07 & CO($2$--$1$) & +172.3 & 2 & 310--455 & 6.4 & -0.16 & 4.2 & unattributed; no recurrence \\
HD 69830 & X12209f7\_X16edb & 7 & 322.4325 & 322.4444 & 5.259 & 5.42 & C$^{18}$O($3$--$2$) & -6268.5 & 5 & 0.1--29 & 8.9 & -0.48 & 4.8 & unattributed; no recurrence \\
ALMA J153702653-33192492 & Xfe62c1\_X398 & 6 & 216.1844 & 216.2036 & 5.228 & 6.02 & SiO($5$--$4$) & -1244.6 & 7 & 18--126 & 11.0 & +1.03 & 4.3 & unattributed; no recurrence \\
AU Mic & X899169\_X1838 & 6 & 230.8252 & 230.8294 & 5.221 & 5.85 & CO($2$--$1$) & +379.0 & 2 & 145--310 & 10.0 & +0.19 & 4.5 & unattributed; no recurrence \\
HD 14055 & Xff99e1\_X2c3bb & 7 & 345.5256 & 345.5216 & 5.186 & 5.66 & SO($2_{3}$--$2_{1}$) & -158.6 & 30 & 0.1--651 & 172.8 & +2.00 & 5.1 & unattributed; no recurrence \\
HD 14055 & Xfffde1\_X3835 & 7 & 345.0099 & 345.0070 & 5.177 & 6.56 & SO($2_{3}$--$2_{1}$) & -604.8 & 30 & 0.1--649 & 187.2 & -1.78 & 5.2 & unattributed; no recurrence \\
HD 14055 & Xff99e1\_X2b2ee & 7 & 344.7007 & 344.6965 & 5.175 & 5.92 & SO($8_{8}$--$7_{7}$) & +336.0 & 30 & 0.1--651 & 210.7 & -0.00 & 5.2 & unattributed; no recurrence \\
LHS 1140 & Xd9d7c7\_X6fd2 & 6 & 230.6355 & 230.6268 & 5.136 & 5.68 & CO($2$--$1$) & +115.5 & 4 & 5.0--27 & 7.3 & +0.33 & 4.0 & unattributed; no recurrence \\
HD 10647 & Xff294f\_Xaf70 & 7 & 331.6570 & 331.6943 & 5.120 & 5.91 & $^{13}$CO($3$--$2$) & +1003.3 & 10 & 0.1--592 & 125.8 & +0.39 & 5.1 & unattributed; no recurrence \\
CD$-$57 1054 & Xcf525f\_X5aa4 & 7 & 346.5099 & 346.5265 & 5.108 & 5.71 & SO($9_{8}$--$8_{7}$) & -1.7$^{\dagger}$ & 1 & 39 & 6.0 & -0.27 & 4.1 & unattributed; no recurrence \\
ALMA J153702653-33192492 & Xf8f6a9\_X10359 & 6 & 217.7292 & 217.7269 & 5.091 & 5.41 & H$_2$CO($3_{0,3}$--$2_{0,2}$) & -680.5 & 7 & 3.8--242 & 11.6 & +0.07 & 4.6 & unattributed; no recurrence \\
GJ 849 & Xd9668b\_X841a & 6 & 230.8769 & 230.8581 & 5.090 & 5.83 & CO($2$--$1$) & +416.3 & 5 & 0.1--51 & 9.3 & -1.69 & 5.3 & unattributed; no recurrence \\
ALMA J153702653-33192492 & Xfe62c1\_X398 & 6 & 218.8729 & 218.8923 & 5.090 & 5.56 & H$_2$CO($3_{2,1}$--$2_{2,0}$) & +181.2 & 7 & 18--126 & 10.5 & -1.04 & 5.0 & unattributed; no recurrence \\
HD 14055 & X1001fdc\_X29ae & 7 & 331.1559 & 331.1546 & 5.082 & 6.26 & $^{13}$CO($3$--$2$) & +513.8 & --- & --- & --- & --- & --- & unattributed; no spectrum \\
$\beta$~Pic & Xd9668b\_Xa9df & 6 & 241.8367 & 241.8594 & 5.079 & 5.99 & CS($5$--$4$) & -3765.1 & 1 & 1.0 & 4.8 & -0.63 & 3.9 & unattributed; no recurrence \\
ALMA J153702653-33192492 & Xf8f6a9\_X10359 & 6 & 218.6711 & 218.6688 & 5.048 & 5.77 & H$_2$CO($3_{2,1}$--$2_{2,0}$) & -125.1 & 7 & 3.8--242 & 11.6 & +0.00 & 4.6 & unattributed; no recurrence \\
HD 69830 & X12277ed\_X178c6 & 7 & 321.2005 & 321.2149 & 5.046 & 6.13 & C$^{18}$O($3$--$2$) & -7387.7 & 5 & 6.0--36 & 10.2 & -0.46 & 4.7 & unattributed; no recurrence \\
HR 1010 & Xc66ce1\_Xaa9 & 6 & 230.8348 & 230.8491 & 5.044 & 5.72 & CO($2$--$1$) & +404.6 & 7 & 1.1--14 & 10.3 & -1.53 & 5.2 & unattributed; no recurrence \\
61 Vir & Xc067f7\_X822 & 7 & 345.5599 & 345.5678 & 5.038 & 5.79 & SO($2_{3}$--$2_{1}$) & -118.5 & 3 & 2.0--3.0 & 9.3 & +0.45 & 4.2 & unattributed; no recurrence \\
HD 170773 & Xff0235\_X8c42 & 7 & 331.0101 & 331.0234 & 5.029 & 5.75 & $^{13}$CO($3$--$2$) & +394.8 & 4 & 0.9--113 & 74.8 & +1.56 & 4.9 & unattributed; no recurrence \\
ALMA J153702653-33192492 & Xd9d7c7\_Xbb64 & 7 & 347.2179 & 347.1890 & 5.009 & 5.68 & SiO($8$--$7$) & -122.2 & 1 & 0.1 & 5.0 & -2.06 & 5.0 & unattributed; no recurrence \\
HD~14055 & X1190de8\_Xd526 & 7 & 344.4676 & 344.4515 & 5.006 & 5.96 & SO($8_{8}$--$7_{7}$) & +122.7 & 30 & 20--611 & 11.6 & -0.60 & 4.8 & unattributed; no recurrence \\
$\eta$~Crv & X122b6ff\_X1041e & 7 & 357.2815 & 357.2475 & 5.006 & 5.50 & HCO$^{+}$($4$--$3$) & +431.3 & 1 & 87 & 11.0 & -0.62 & 4.8 & unattributed; no recurrence \\
\hline
\end{tabular}

\end{table*}
